\subsection{Comparativa de herramientas GRC y de apoyo al cumplimiento}\label{comparativa-de-herramientas-grc-y-de-apoyo-al-cumplimiento}

\subsection{Universo evaluado}\label{universo-evaluado}

Se compararon \textbf{ocho plataformas}: las seis de la ficha ---OneTrust, Vanta, Drata, BigID, Microsoft Purview Compliance Manager y Securiti.ai--- más dos de aporte propio, \textbf{Osano} y \textbf{Eramba}. Osano cubre el tramo de la organización mediana que necesita gestión de consentimiento y no una suite completa, que ninguna de las seis atiende. Eramba introduce la única alternativa \textbf{autoalojada} del conjunto, lo que permite incorporar el eje construir frente a comprar en vez de limitar la comparación a proveedores SaaS; sin ella el universo habría tenido un sesgo de categoría.

Cada una se caracterizó con cinco columnas objetivas ---cobertura de la Ley 21.719 con plantilla propia, cobertura ISO 27001/27701/42001, descubrimiento de datos, gestión de consentimiento y publicación de precio--- tomadas de su página oficial el 20-09-2026. \textbf{El dato decisivo es que solo dos de las ocho publican precio}: Purview por plan de licenciamiento y Eramba por edición comunitaria. Las seis restantes operan exclusivamente por cotización.

\subsection{Criterios, ponderación y análisis}\label{criterios-ponderaciuxf3n-y-anuxe1lisis}

\subsection{Criterios ponderados (los que sí distinguen)}\label{criterios-ponderados-los-que-suxed-distinguen}

{\def\LTcaptype{none} % do not increment counter
\begin{longtable}[]{@{}
  >{\raggedright\arraybackslash}p{(\linewidth - 6\tabcolsep) * \real{0.2353}}
  >{\centering\arraybackslash}p{(\linewidth - 6\tabcolsep) * \real{0.2941}}
  >{\raggedright\arraybackslash}p{(\linewidth - 6\tabcolsep) * \real{0.2353}}
  >{\raggedright\arraybackslash}p{(\linewidth - 6\tabcolsep) * \real{0.2353}}@{}}
\toprule\noalign{}
\begin{minipage}[b]{\linewidth}\raggedright
Criterio
\end{minipage} & \begin{minipage}[b]{\linewidth}\centering
Peso
\end{minipage} & \begin{minipage}[b]{\linewidth}\raggedright
Por qué pesa lo que pesa
\end{minipage} & \begin{minipage}[b]{\linewidth}\raggedright
De dónde sale el dato
\end{minipage} \\
\midrule\noalign{}
\endhead
\bottomrule\noalign{}
\endlastfoot
Descubrimiento de datos personales & 25 \% & Curimón no tiene inventario de tratamientos. Sin esto, el registro de actividades se levanta a mano sobre la base operativa completa & Columna «Descubrimiento de datos» \\
Gestión de consentimiento & 20 \% & 258 conductores externos y 148 transportistas generan solicitudes de acceso y supresión que alguien debe resolver en plazo & Columna «Consentimiento» \\
Transparencia del precio & 25 \% & El punto 5 de las Indicaciones impide llevar al flujo de caja una cifra no verificable. Solo 2 de 8 publican precio real & Columna «Precio publicado» \\
Encaje con la arquitectura comprometida & 20 \% & La propuesta ya fijó Azure y Key Vault con HSM en el Informe 1. Una herramienta nativa del mismo ecosistema reduce costo de integración & Modelo de despliegue y región declarada por el proveedor \\
Esfuerzo de operación & 10 \% & La única autoalojada del conjunto es Eramba. Ahí el esfuerzo de operación es el criterio que decide si conviene frente a comprar SaaS & Inferido del modelo de despliegue; no figura como columna en la ficha del §3.2 \\
\end{longtable}
}

Suma: 100 \%. La ponderación es una decisión del equipo y admite revisión: si se asume que toda contratación pasará igualmente por cotización, el 25 \% de transparencia del precio pierde poder discriminante y se traslada al encaje con la arquitectura, que es el criterio de mayor efecto económico en este caso.

\subsection{Matriz de puntuación}\label{matriz-de-puntuaciuxf3n}

Escala 1 a 5. Regla de conversión: Sí = 5, Parcial o «Demo» = 3, No = 1. Precio: publicado = 5, solo cotización = 1 (sin matices, porque el punto 5 trata ambos casos igual: si no hay precio de lista, no hay precio de lista).

{\def\LTcaptype{none} % do not increment counter
\begin{longtable}[]{@{}
  >{\raggedright\arraybackslash}p{(\linewidth - 12\tabcolsep) * \real{0.1818}}
  >{\centering\arraybackslash}p{(\linewidth - 12\tabcolsep) * \real{0.1364}}
  >{\centering\arraybackslash}p{(\linewidth - 12\tabcolsep) * \real{0.1364}}
  >{\centering\arraybackslash}p{(\linewidth - 12\tabcolsep) * \real{0.1364}}
  >{\centering\arraybackslash}p{(\linewidth - 12\tabcolsep) * \real{0.1364}}
  >{\centering\arraybackslash}p{(\linewidth - 12\tabcolsep) * \real{0.1364}}
  >{\centering\arraybackslash}p{(\linewidth - 12\tabcolsep) * \real{0.1364}}@{}}
\toprule\noalign{}
\begin{minipage}[b]{\linewidth}\raggedright
Herramienta
\end{minipage} & \begin{minipage}[b]{\linewidth}\centering
Descubr. 25 \%
\end{minipage} & \begin{minipage}[b]{\linewidth}\centering
Consent. 20 \%
\end{minipage} & \begin{minipage}[b]{\linewidth}\centering
Precio 25 \%
\end{minipage} & \begin{minipage}[b]{\linewidth}\centering
Encaje Azure 20 \%
\end{minipage} & \begin{minipage}[b]{\linewidth}\centering
Operación 10 \%
\end{minipage} & \begin{minipage}[b]{\linewidth}\centering
\textbf{Total ponderado}
\end{minipage} \\
\midrule\noalign{}
\endhead
\bottomrule\noalign{}
\endlastfoot
OneTrust & 5 & 5 & 1 & 1 & 3 & \textbf{3,05} \\
Vanta & 3 & 1 & 1 & 1 & 3 & \textbf{1,75} \\
Drata & 3 & 1 & 1 & 1 & 3 & \textbf{1,75} \\
BigID & 5 & 3 & 1 & 1 & 3 & \textbf{2,55} \\
\textbf{Microsoft Purview} & 5 & 1 & \textbf{5} & \textbf{5} & 4 & \textbf{3,70} \\
Securiti.ai & 5 & 5 & 1 & 1 & 3 & \textbf{3,05} \\
Osano & 3 & 5 & 1 & 1 & 4 & \textbf{2,65} \\
Eramba & 1 & 1 & \textbf{5} & 1 & 1 & \textbf{1,60} \\
\end{longtable}
}

\subsection{Lectura de la matriz y recomendación}\label{lectura-de-la-matriz-y-recomendaciuxf3n}

La mejor evaluada es \textbf{Microsoft Purview, con 3,70 puntos}, y gana por una vía que conviene explicitar porque no es la esperable: \textbf{no es la más completa funcionalmente}. Puntúa 1 en gestión de consentimiento, el peor valor posible en un criterio que pesa 20 \%, y aun así se impone, porque es la única que obtiene el máximo simultáneo en los dos criterios que concentran el 45 \% de la ponderación: transparencia del precio y encaje con la arquitectura comprometida. El segundo lugar lo comparten OneTrust y Securiti.ai con 3,05; ambas son funcionalmente superiores ---5 en descubrimiento y 5 en consentimiento--- y pierden porque ninguna publica precio ni es nativa del ecosistema ya fijado. \textbf{La distancia entre 3,70 y 3,05 no mide capacidad técnica: mide verificabilidad y costo de integración.} Eramba ilustra el límite por el otro extremo: única con precio íntegramente público, obtiene el peor total (1,60) porque puntúa el mínimo en descubrimiento, consentimiento y encaje, de modo que el ahorro de licencia se traslada íntegro a esfuerzo de operación propio.

La propuesta, sin embargo, \textbf{adopta CISO Assistant Pro}, y la razón está en el criterio que la matriz pondera más alto. Purview publica \textbf{por plan de licenciamiento}, no una tarifa autónoma: su partida dependería del plan Microsoft 365 o Azure que se contrate, y ese plan no está fijado. CISO Assistant Pro tiene \textbf{precio anual firme y público por ser software libre AGPLv3}, lo que permite presupuestarlo sin declarar un supuesto, y el punto 5 de las Indicaciones es categórico en que una cifra que no puede verificarse no entra en una estimación.

Dos consecuencias se declaran con la misma franqueza. La primera: \textbf{CISO Assistant Pro no formó parte del universo de ocho comparadas}; entró por la vía del modelo económico al constatarse que seis de las ocho operan solo por cotización, y no se le asigna puntuación de una matriz que no corrió. La segunda: \textbf{ninguna de las dos resuelve la gestión de consentimiento}, que el caso necesita para los 258 conductores externos y los 148 transportistas sin vínculo laboral; ese tramo se resuelve en el módulo de consentimiento de la app móvil del conductor, no por la vía de la herramienta GRC.
